\documentclass[11pt]{article}
\usepackage{ochamath}
\subjectcolor{2F5DA8}
\setsheet{線形代数}{線形写像と基底変換　演習}{https://university-math-crj.pages.dev/linear-algebra/linear-maps/}
\namelinetrue
\begin{document}
\makesheethead{問題}

\problem[線形性]
$T(x,y)=(2x-y,3y)$ と $S(x,y)=(x+1,y)$ は線形か。理由を答えよ。
\workspace{38mm}

\problem[核と像]
$T(x,y,z)=(x+z,y+z)$ の核と像の基底を求め、次元を確認せよ。
\workspace{38mm}
\newpage

\problem[同じ基底での表現]
$A=\begin{pmatrix}2&3\\0&5\end{pmatrix}$、$P=\begin{pmatrix}1&1\\0&1\end{pmatrix}$ とする。$P^{-1}AP$ を求めよ。
\workspace{38mm}

\problem[異なる基底]
標準座標で恒等写像 $A=I_2$ を、入力基底 $P=\operatorname{diag}(2,1)$、出力基底 $Q=\operatorname{diag}(1,3)$ で表せ。入力基底座標 $(1,3)^{\mathsf T}$ の出力座標も求めよ。
\workspace{38mm}
\end{document}
